\section{Discussion and open questions}\label{sec:discussion}

\subsection{The first-order rule is no guide}
The first-order rules $54$ and $110$ are the elementary rules associated with class~4 behaviour \cite[pp.~693--697]{Wolfram2002}; their second-order versions have automorphism groups of opposite kinds, and $110$R and $30$R, whose first-order rules behave differently, have the same group $\Z^2$ and the same commuting block maps through radius $4$. What distinguishes $54$R is a conserved local structure, the wall $33$, with an autonomous interior. The same mechanism, in the forms of fixed, hidden-bit, time-periodic and rotating walls, accounts for $39$ of the $43$ classes of status (G); the walls of these classes go back to Takesue \cite[Section~4]{Takesue1989}, and their use for automorphisms is that of \cite{PaperC}. Three of the other four, $90$, $165$ and $255$, are affine \cite{PaperC}. The fourth is the class of $95$, which is not affine: there the room comes from the particle-conserving automorphisms of a gas of light-speed particles, and $F_{95}$ has no blocking word and no eigenfield of any speed \eqref{imp:95}. The particles of the first-order rule $54$ are well studied \cite{BNR1991,MACC2012}; they play no role here. The reversible automaton called Rule~54 in \cite{BBGP1993,BKP2021} is neither rule $54$ nor $54$R: it is the system on one sublattice of a spacetime diagram of the class of $95$ \cite{PaperE}.

\subsection{Rigidity}
On the open classes the edge-move argument of \cite{PaperA} does not get started in $21$ classes, for lack of a stuck edge or of a kink next to a background with one, and stops later in the computations for $6$ more (Section~\ref{cl:sec:reach}). For $91$ Lemma~L holds at every radius, but no normalisation of the automorphisms is known; for $155$ even the proof of Lemma~L fails \eqref{imp:L}. For the $21$ open classes with finite $\Phi_f$ there are no kinks of light speed from which to read a homomorphism to $\Z^2$.

Proposition~\ref{cl:prop:nowalls} shows that $110$R has no fixed, time-periodic, rotating or hidden-bit wall of any length; the rigidity of $110$R \cite{PaperA} is the stronger statement that it has no automorphism outside $\langle\sh,F\rangle$ of any kind. Conjugacy of reversible automata is undecidable \cite{Salo2020}; we do not know whether the property $\Aut(X_F)=\langle\sh,F\rangle$ of a reversible automaton $F$ is decidable.

\subsection{Open questions}
A word fixed in every context is blocking (Section~\ref{sec:blocking}), and the walls behind the status (G) entries other than $90$, $95$, $165$ and $255$ all give blocking words \cite{PaperC}; for $255$ every word is blocking. Among the $43$ classes of status (G), those of $90$, $165$ and $95$ are the only ones without a blocking word \eqref{imp:90}, \eqref{imp:95}. Question~1.9 of \cite{PaperC} asks whether, for an elementary rule $f$ that is not affine and does not belong to the class of $95$, the absence of a blocking word for $F_f$ implies that $\Aut(X_{F_f})$ is virtually $\Z^2$. The classification answers it for the classes of status (G), (Z), (S) and (A): the classes of status (A) are affine, the classes of status (G) that are not affine, other than that of $95$, have blocking words, and in the classes of status (Z) and (S) the group is virtually $\Z^2$. So that question concerns the $28$ open classes.

\begin{question}\label{q:open}
Is $\Aut(X_{F_f})=\langle\sh,F_f\rangle$ for each of the $28$ open classes?
\end{question}

The evidence is weak: the class of $95$ passes the same censuses. The natural candidate for a second exception is the class of $175$, the only open class whose rule does not read its middle argument, so that the two sublattices of a diagram evolve independently, as for $95$ \cite[Remark~6.22]{PaperE}.

\begin{enumerate}
\item Decide the $28$ open classes. Possible routes: for the class of $25$, the remaining open class whose light-speed fixed points include kinks as for $110$, and for the class of $31$, where kinks of $\sh^{\pm1}F^3$ and moves behave as for $110$ in the computations of Remark~\ref{cl:rem:reach}, an argument as in \cite{PaperA}; for $91$, a normalisation of the automorphisms as for $94$, and for $155$, a proof of Lemma~L (see the questions in \cite{PaperE}). For the $21$ open classes with finite $\Phi_f$ a homomorphism $\Aut(X_{F_f})\to\Z^2$ extending $\sh^aF^b\mapsto(a,b)$ is missing.
\item Extend the census beyond radius $4$ for the open classes.
\item The following are questions of \cite{PaperC}. Determine $\Aut(X_{F_f})$ for a class of status (G) with walls. For the affine classes of statuses (S) and (A), is every block map commuting with $F_f$ affine? Every automorphism is \eqref{imp:affine}, and every commuting block map of radius at most $3$ is (computer-assisted in \cite{PaperC}). For $54$R, is $\Aut(X_{F_{54}})$ generated by $\sh$, $F$ and the wall relabellings and wall-trading involutions of the kind in \cite{PaperC}, of all radii? For $90$R, is it generated by the affine automorphisms and the marker maps of \cite{PaperC}?
\item For the class of $95$, describe the group $\mathcal N$ of \eqref{imp:95} and the block maps that commute with $F_{95}$ and are not invertible; see the questions in \cite{PaperE}.
\end{enumerate}
